% Chapter - Modelling, dependency and correlation ............................................................

\chapter{Modelling, dependency, and correlation}
\label{chapter5}

\epigraph{\textit{Numbers have an important story to tell,\\if given a voice.}}{— Florence Nightingale}

The modern scientific conception of nature rests on the idea that observable phenomena can be represented through variables whose relationships follow general mathematical laws, a perspective made explicit through analytic geometry and the introduction of coordinate systems that allowed geometric relations to be expressed algebraically \cite{descartes_1637_geometry,newton_1687_principia}. 

\medskip

The development of calculus extended this framework by providing tools to describe continuous change, enabling relationships between quantities to be expressed as functions \(y = f(x)\) whose rates of change could be studied through derivatives and whose accumulated effects could be analyzed through integration \cite{leibniz_1684_nova,cauchy_1821_analyse}. 

\medskip

Prediction became the defining feature of scientific explanation: if the relation between variables can be expressed mathematically, then under specified assumptions one may compute future or unobserved values, thereby linking empirical regularities with formal reasoning and laying conceptual foundations later absorbed into probabilistic modelling \cite{reichenbach_1938_prediction,cox_2006_principles}. 

% Section 1 ...................................................................................................

\section{The ide of model, explanation and prediction}

Scientific investigation proceeds from observation to structured explanation, transforming collections of measurements into relations among quantities that can be expressed symbolically and analyzed systematically \cite{hacking_1975_emergence}. 

\medskip

A fundamental conceptual distinction in modelling is between \emph{variables} and \emph{parameters}. A variable is a quantity that may take different values across individuals, experiments, or time, whereas a parameter is a fixed but unknown constant that characterizes the underlying mechanism generating those variable outcomes. 

\medskip

The distinction emerged gradually in early probability theory: Bernoulli treated probabilities as fixed characteristics of repeatable processes \cite{bernoulli_1713_ars}, Laplace formalized parameters within analytical probability models \cite{laplace_1812_theorie}, and in the twentieth century Fisher clarified parameters as fixed constants to be estimated from random samples \cite{fisher_1925_statistical}. 

\medskip

Formally, if \(Y_1,\dots,Y_n\) denote random variables representing observed measurements, one may posit that their distribution depends on a parameter \(\theta\), written \(Y_i \sim f(y \mid \theta)\), where \(\theta\) is unknown but not random in the classical frequentist framework. 

\medskip

Thus in a regression model such as 
\[
Y = a + bX + \varepsilon,
\]
the quantities \(X\) and \(Y\) are variables that vary across observations, while \(a\) and \(b\) are parameters describing the systematic structure relating them. 

\medskip

For laboratory scientists, this distinction is crucial: measurements such as enzyme activity, concentration, or response time are variables that fluctuate, whereas quantities such as a true mean effect or regression slope are parameters summarizing the underlying biological mechanism that generated those measurements \cite{degroot_2012_probability}.

% Section 2 ...................................................................................................

\section{Probability Structure, Error, and Uncertainty}

All empirical measurement exhibits variability, and Gauss’s theory of errors formalized the distinction between systematic structure and random deviation in scientific observation \cite{gauss_1809_theoria}. 

\medskip

Kolmogorov’s axiomatization of probability \cite{kolmogorov_1933_grundbegriffe} provided the mathematical foundation for expressing models in the form \(Y = f(X) + \varepsilon\), separating predictable dependence from stochastic fluctuation. 

\medskip

Twentieth-century developments reframed error not as failure but as intrinsic uncertainty requiring quantification rather than elimination \cite{cox_2006_principles,mayo_1996_error}. 

\medskip

For experimental researchers, the error term represents biological variability, measurement imprecision, and unobserved influences, and properly accounting for it guards against overstated conclusions and irreproducible findings \cite{ioannidis_2005_false}. 

% Section 3 ...................................................................................................

\section{Estimation, Inference, and Scientific Evidence}

Once a model specifies variables and parameters, the central inferential task is to estimate the unknown parameters from observed data. Fisher systematized this process through likelihood theory, treating parameters as fixed but unknown constants and the observed data as realizations of random variables \cite{fisher_1925_statistical}. 

\medskip

Given independent observations \(Y_1,\dots,Y_n\) with density \(f(y \mid \theta)\), the likelihood function is defined as 
\[
L(\theta) = \prod_{i=1}^n f(Y_i \mid \theta),
\]
and the maximum likelihood estimator is 
\[
\hat{\theta} = \arg\max_{\theta} L(\theta).
\]

\medskip

This framework formalizes the conceptual separation between random variables and fixed parameters: variability arises from the sampling distribution of the estimator \(\hat{\theta}\), not from randomness in \(\theta\) itself within the classical paradigm \cite{degroot_2012_probability}. 

\medskip

Neyman and Pearson further developed procedures for testing hypotheses about parameters while controlling long-run error rates \cite{pearson_neyman_1933_efficient}, establishing systematic rules for deciding whether observed data provide sufficient evidence to reject a specified parameter value. 

\medskip

For laboratory scientists, this means that while individual measurements vary unpredictably, the target of inference is a stable underlying quantity — such as a true mean difference, rate constant, or regression slope — whose uncertainty is quantified through confidence intervals and hypothesis tests grounded in sampling theory \cite{cox_2006_principles}.

% Section 4 ...................................................................................................

\section{Linear Dependence, Correlation, and Linear Models}

The concept of regression originated in Galton’s studies of heredity and was formalized mathematically by Pearson through correlation and least-squares methods \cite{pearson_1895_contributions,spearman_1904_association}. 

\medskip

A simple linear regression model takes the form \(Y = a + bX + \varepsilon\), where the slope parameter \(b\) represents the expected change in \(Y\) associated with a one-unit change in \(X\). 

\medskip

This framework generalizes to the linear model \(Y = X\beta + \varepsilon\), unifying regression, analysis of variance, and experimental design within a single structure \cite{fisher_1925_statistical,fisher_1935_design}. 

\medskip

For laboratory researchers, linear models allow simultaneous adjustment for multiple predictors such as treatment, age, and batch effects, while correlation provides a summary of association without implying causation. 

% Section 5 ...................................................................................................

\section{Beyond Linearity: Generalized Linear Models and Interpretation}

Many scientific outcomes are not normally distributed, and the probabilistic foundation established by Kolmogorov \cite{kolmogorov_1933_grundbegriffe} permits systematic extension beyond classical linear assumptions. 

\medskip

Generalized Linear Models retain a linear predictor while allowing alternative response distributions through a link function, expressed abstractly as \(g(\mathbb{E}[Y]) = X\beta\). 

\medskip

Logistic regression models binary outcomes, Poisson regression models counts, and modern biological analyses such as RNA-seq differential expression rely on likelihood-based GLM frameworks \cite{love_2014_deseq2,cox_2006_principles}. 

\medskip

Finally, statistical association does not automatically establish causation, and rigorous interpretation requires careful design, testing, and critical scrutiny consistent with principles of scientific reasoning \cite{popper_1934_logic,mayo_2018_severe}. 
